\documentclass[11pt]{article}

% ---------- Page ----------
\usepackage[margin=1in]{geometry}

% ---------- Engine-safe fonts (XeLaTeX / LuaLaTeX) ----------
\usepackage{fontspec}
\setmainfont{Latin Modern Roman}
\setsansfont{Latin Modern Sans}
\setmonofont{Inconsolata}

% ---------- Typography ----------
\usepackage{microtype}
\usepackage{setspace}
\setstretch{1.08}

% ---------- Links ----------
\usepackage[hidelinks]{hyperref}

% ---------- Colors ----------
\usepackage{xcolor}
\definecolor{ink}{HTML}{111827}
\definecolor{sub}{HTML}{374151}
\definecolor{muted}{HTML}{6B7280}
\definecolor{line}{HTML}{E5E7EB}
\definecolor{panel}{HTML}{F9FAFB}
\definecolor{accent}{HTML}{2563EB}
\definecolor{codebg}{HTML}{111827}

% ---------- Headings ----------
\usepackage{titlesec}
\titleformat{\section}
  {\sffamily\Large\bfseries\color{ink}}{\thesection}{0.75em}{}
\titleformat{\subsection}
  {\sffamily\large\bfseries\color{ink}}{\thesubsection}{0.75em}{}
\titleformat{\subsubsection}
  {\sffamily\normalsize\bfseries\color{ink}}{\thesubsubsection}{0.75em}{}

% ---------- Lists ----------
\usepackage{enumitem}
\setlist[itemize]{topsep=4pt,itemsep=2pt,leftmargin=*}
\setlist[enumerate]{topsep=4pt,itemsep=2pt,leftmargin=*}

% ---------- Modern boxes ----------
\usepackage[most]{tcolorbox}
\tcbset{
  enhanced,
  boxrule=0.6pt,
  colframe=line,
  colback=panel,
  arc=8pt,
  left=10pt,right=10pt,top=8pt,bottom=8pt
}

% ---------- Code (minted) ----------
\usepackage{minted}
\setminted{
  style=monokai,
  bgcolor=codebg,
  fontsize=\small,
  linenos,
  breaklines,
  autogobble
}

% ---------- Custom code/output boxes ----------
\newtcolorbox{cmd}[1][]{%
  title={\sffamily\bfseries Command%
    \if\relax\detokenize{#1}\relax\else\ — #1\fi},
  coltitle=ink,
  colback=panel,
  colframe=line
}

\newtcolorbox{out}[1][]{%
  title={\sffamily\bfseries Output%
    \if\relax\detokenize{#1}\relax\else\ — #1\fi},
  coltitle=ink,
  colback=panel,
  colframe=line
}

% ---------- Header / Footer ----------
\usepackage{fancyhdr}
\pagestyle{fancy}
\fancyhf{}
\setlength{\headheight}{14pt}
\lhead{\textcolor{muted}{MongoDB Indexing Lab}}
\rhead{\textcolor{muted}{\today}}
\cfoot{\textcolor{muted}{\thepage}}

% ---------- Title ----------
\usepackage{titling}
\pretitle{\vspace{-1.0cm}\color{ink}\LARGE\sffamily\bfseries}
\posttitle{\par\vspace{0.15cm}\color{sub}\large\sffamily}
\preauthor{\color{sub}\normalsize\sffamily}
\postauthor{\par}
\predate{\color{muted}\normalsize\sffamily}
\postdate{\par\vspace{0.6cm}\hrule\vspace{0.8cm}}

\title{Indexing and Query Performance in MongoDB}
\author{Nuria Arroyo \quad \texttt{(173605)}}
\date{MongoDB Compass + Atlas \; \textbullet \; Database \texttt{tp1}, Collection \texttt{restaurantes}}

\begin{document}
\maketitle

\begin{tcolorbox}
\textbf{Objective.} Explore the impact of indexing on query performance in MongoDB. We run baseline queries without indexes, inspect execution plans using \texttt{explain}, create the required indexes, re-run the same queries, and compare performance metrics. Finally, we drop the indexes and observe the regression.
\end{tcolorbox}

\section{Environment}
\begin{itemize}
  \item Platform: MongoDB Atlas + MongoDB Compass
  \item Database: \texttt{tp1}
  \item Collection: \texttt{restaurantes}
  \item Dataset size (approx.): \underline{\hspace{3cm}} documents
\end{itemize}

\section{Procedure (Assignment Instructions)}
\begin{enumerate}
  \item Execute queries \textbf{without indexes} and observe performance using \texttt{explain}.
  \item Create the specified indexes (Section 3.5).
  \item View / verify indexes.
  \item Execute the same queries \textbf{with indexes} and observe changes.
  \item Compare and analyze results.
  \item Delete indexes and re-check query performance.
\end{enumerate}

\newpage

\section{Step 1 — Baseline: queries without indexes}
\subsection{Confirm existing indexes}
\begin{cmd}[Verify index state (baseline)]
\begin{minted}{javascript}
use tp1
db.restaurantes.getIndexes()
\end{minted}
\end{cmd}

\begin{out}[Paste your output here]
\begin{minted}[linenos=false]{text}
% Paste output here
\end{minted}
\end{out}

\subsection{Baseline query: all Italian restaurants}
\begin{cmd}[Explain with execution stats (baseline)]
\begin{minted}{javascript}
db.restaurantes.find({ cuisine: "Italian" }).explain("executionStats")
\end{minted}
\end{cmd}

\begin{out}[Paste your explain output here]
\begin{minted}[linenos=false]{text}
% Paste explain output here
\end{minted}
\end{out}

\section{Step 2 — Create indexes (Section 3.5)}
\begin{cmd}[Create indexes]
\begin{minted}{javascript}
db.restaurantes.createIndex({ cuisine: 1 })
db.restaurantes.createIndex({ cuisine: 1, "address.zipcode": -1 })
\end{minted}
\end{cmd}

\section{Step 6 — Delete indexes and re-check}
\begin{cmd}[Drop indexes]
\begin{minted}{javascript}
db.restaurantes.dropIndex("cuisine_1")
db.restaurantes.dropIndex("cuisine_1_address.zipcode_-1")
\end{minted}
\end{cmd}

\end{document}